
\chapter{Guards, Queries и выбор перехода}
\label{chap:guards-queries}

До сих пор мы рассматривали переходы, которые
определялись текущим состоянием автомата и поступившей
командой. Однако иногда для одной команды необходимо
предусмотреть несколько вариантов поведения.

Для этого в CIAO используются условия Guard.
Они позволяют выбирать переход на основе данных
контекста или результатов запросов к другим автоматам.

\section{Условные переходы}

Рассмотрим упрощённый фрагмент Array Sum.
При получении команды \code{tick} автомат должен
проверить, остались ли элементы для обработки.

\begin{lstlisting}[language=Python,numbers=none]
@guard(text="есть следующий элемент")
def has_next(self, ctx: SumContext) -> bool:
    return ctx.index < len(ctx.values)

on_tick = transition(
    source=running,
    event=tick,
    guard=has_next,
    actions=(accumulate,),
    target=running,
)

on_complete = transition(
    source=running,
    event=tick,
    otherwise=True,
    actions=(complete,),
    target=publishing,
)
\end{lstlisting}

Метод \code{has_next} объявлен с помощью
\code{@guard} и возвращает \code{bool}.
Его результат определяет, можно ли выполнить
переход \code{on_tick}.

Если элементов больше нет, выбирается резервная
ветвь \code{on_complete}, обозначенная
\code{otherwise=True}. Она выполняется, только
когда ни один Guard группы не вернул \code{True}.

На диаграмме это выглядит следующим образом:

\noindent\begin{minipage}{\textwidth}
  \centering
  \begin{tikzpicture}[
  state/.style={
    draw=ciaoBlue,
    fill=ciaoPaleBlue!45,
    very thick,
    rounded corners=2mm,
    minimum width=36mm,
    minimum height=16mm,
    font=\sffamily\bfseries
  },
  flow/.style={-{Stealth[length=3mm]},thick,draw=ciaoNavy},
  transition label/.style={
    fill=white,
    inner sep=2pt,
    align=center,
    font=\ttfamily\small
  }
]
  \node[state] (running) {running};
  \node[state,right=55mm of running] (publishing) {publishing};

  % Guard-ветвь остаётся в running, поэтому показываем её self-loop.
  \draw[flow]
    (running) to[out=120,in=60,looseness=5]
    node[transition label,above=1mm]
      {[has\_next: есть следующий элемент]\\[-0.3mm]
       / accumulate; tick}
    (running);

  % Резервная ветвь выбирается только при ложных Guard группы.
  \draw[flow]
    (running) --
    node[transition label,below=2mm]
      {[else] / complete; tick}
    (publishing);
\end{tikzpicture}

  \captionof{figure}{Guard и резервная ветвь группы переходов}
  \label{fig:guard-transition}
\end{minipage}

\subsection{Правила выбора перехода}

Переходы с одинаковыми \code{source} и \code{event}
образуют одну группу. При обработке команды Runtime
проверяет все Guard этой группы в порядке объявления.

Возможны следующие результаты:

\begin{itemize}
  \item Один Guard вернул \code{True} ---
    выбирается соответствующий переход.
  \item Несколько Guard вернули \code{True} ---
    возникает \code{AmbiguousTransitionError}.
  \item Все Guard вернули \code{False} ---
    выбирается ветвь \code{otherwise}, если она есть.
\end{itemize}

Таким образом, порядок объявления переходов
не задаёт их приоритет. Условия должны исключать
одновременный выбор нескольких ветвей.

Если условная группа не содержит \code{otherwise},
компилятор модели CIAO выдаёт предупреждение
\code{MissingOtherwiseWarning}. Когда все условия
ложны, Runtime обрабатывает ситуацию как отсутствие
подходящего перехода.

\begin{importantbox}
Guard должен только проверять условие, не изменяя
контекст автомата. Изменения следует выполнять
в Action выбранного перехода.
\end{importantbox}

\section{Синхронные запросы}

Иногда для проверки условия или выполнения действия
автомату нужны данные другого объекта.

Например, при обработке математического выражения
калькулятору может потребоваться узнать приоритет
операции.

Для таких обращений предусмотрены Query.
В отличие от команд, они выполняются сразу
и возвращают результат в текущий обработчик.

\begin{lstlisting}[language=Python,numbers=none]
@required_query
def precedence(self, symbol: str) -> int: ...

@action(query=precedence)
def inspect(self, ctx: Context, query: Query) -> None:
    priority = query(symbol=ctx.current_symbol)
\end{lstlisting}

Операция \code{precedence} объявлена как
\code{required_query}. Она задаёт параметры запроса
и тип возвращаемого значения.

Декоратор \code{@action(query=precedence)}
указывает, какой запрос используется действием.
При этом в сигнатуре обработчика обязателен
параметр типа \code{Query}.

Теперь объявим предоставляемую операцию в другом
автомате --- \code{OperatorTable}:

\begin{lstlisting}[language=Python,numbers=none]
@provided_query
def precedence(self, symbol: str) -> int:
    priorities = {
        "+": 1, "-": 1,
        "*": 2, "/": 2,
    }
    return priorities.get(symbol, 0)
\end{lstlisting}

В отличие от \code{required_query}, предоставляемый
запрос содержит реализацию. Здесь он возвращает
приоритет переданного арифметического оператора.

Для взаимодействия операции связываются через
\code{Connection}, как и обычные команды:

\begin{lstlisting}[language=Python,numbers=none]
Connection(
    (calculator, Calculator.precedence),
    (operators, OperatorTable.precedence),
)
\end{lstlisting}

На диаграмме это выглядит следующим образом:

\noindent\begin{minipage}{\textwidth}
  \centering
  \begin{tikzpicture}[
  automaton/.style={
    draw=ciaoBlue,
    fill=ciaoPaleBlue!35,
    very thick,
    rounded corners=2mm,
    minimum width=43mm,
    minimum height=23mm,
    align=center,
    font=\sffamily
  },
  query connector/.style={draw=ciaoQueryRed,thick},
  operation/.style={font=\ttfamily\small,color=ciaoQueryRed}
]
  \node[automaton] (calculator) {
    {\scriptsize\color{ciaoMuted}\textless\textless Automaton\textgreater\textgreater}\\[-1mm]
    {\bfseries Calculator}
  };
  \node[automaton,right=67mm of calculator] (operators) {
    {\scriptsize\color{ciaoMuted}\textless\textless Automaton\textgreater\textgreater}\\[-1mm]
    {\bfseries OperatorTable}
  };

  % Required и provided половины Query сходятся в общем assembly-стыке.
  \coordinate (junction) at ($(calculator.east)!0.5!(operators.west)$);
  \draw[query connector]
    (calculator.east) -- ($(junction)+(-3mm,0)$);
  \draw[query connector]
    ($(junction)+(0,3mm)$)
    arc[start angle=90,end angle=270,radius=3mm];

  \coordinate (ball-center) at ($(junction)+(2.4mm,0)$);
  \draw[query connector]
    (operators.west) -- ($(ball-center)+(2.4mm,0)$);
  \draw[query connector,fill=white]
    (ball-center) circle[radius=2.4mm];

  \node[operation,above=6mm of junction] {precedence};
\end{tikzpicture}

  \captionof{figure}{Синхронное Query-соединение \code{precedence}}
  \label{fig:query-connection}
\end{minipage}

Когда Calculator вызывает \code{query(...)},
Runtime направляет запрос объекту \code{operators}
и возвращает результат его обработчику.

Запросы можно использовать не только в Action,
но и в Guard:

\begin{lstlisting}[language=Python,numbers=none]
@guard(query=precedence)
def check_priority(
    self, ctx: Context, query: Query
) -> bool:
    return query(symbol=ctx.current_symbol) > 0
\end{lstlisting}

Если обработчику требуется несколько разных
запросов, используется \code{queries=(...)}
и параметр типа \code{Queries}.

Конкретная операция выбирается по её объявлению:

\begin{lstlisting}[language=Python,numbers=none]
priority = queries[self.precedence](symbol="+")
\end{lstlisting}

Как и для \code{Emit}, соответствие объявленных
операций и параметров обработчика проверяется
при компиляции модели CIAO.

\subsection{Особенности Query}

Query выполняется синхронно, без добавления
сообщения в FIFO. Принимающий автомат отвечает
обычным Python-методом, объявленным через
\code{@provided_query}.

При этом запрос не запускает переход и не изменяет
состояние принимающего автомата.

Объекты \code{Query} и \code{Queries} доступны только
во время вызова соответствующего Guard или Action.
Сохранять их в контексте для последующего
использования нельзя.

\begin{importantbox}
Query рекомендуется использовать только для чтения
данных. Хотя Runtime не запрещает изменять контекст
внутри \code{provided_query}, такие изменения
не откатываются автоматически при ошибках.

Если требуется изменить состояние другого автомата,
используйте Command и явно объявленный переход.
\end{importantbox}
